\begin{textsample}{POS Dim 3 – am – Score 18.00 – am/am\_inf\_26.txt}  \label{ex:f3_pos_039}
12.
O planejamento e a avaliação na Educação \textbf{Infantil} O trabalho com \textbf{bebês} e \textbf{crianças} na Educação \textbf{Infantil} não pode prescindir de uma prática contínua de planejar e replanejar.
Cada momento da \textbf{rotina}, cada experiência proposta, tendo em vista a aprendizagem e desenvolvimento, cada forma de organização do ambiente, do tempo e dos materiais, tudo deve ser sistematicamente analisado.
À equipe pedagógica cabe a tarefa essencial de conhecer os objetivos, motivos e sentidos da ação educativa.
Para tanto, é preciso planejar, projetar, refletir, analisar e avaliar.
Quanto ao documento de \textbf{registro} oficial deste planejamento, este ocorre em periodicidade definida pela proposta pedagógica de cada instituição de ensino, não devendo ultrapassar 30 dias letivos.
Quando se educa, realiza -se um trabalho intencional e todas as ações intervêm e agem sobre o desenvolvimento \textbf{infantil}.
No trabalho com \textbf{bebês} e \textbf{crianças}, não se deve priorizar apenas o desenvolvimento cognitivo, mas todo o conjunto de saberes vividos pelos campos de experiências e que consideram o corpo, as emoções, as vivências.
No tocante à documentação pedagógica, os \textbf{registros} são documentos que permitem rever e repensar sobre as situações para entender o que os \textbf{bebês} e \textbf{crianças} \textbf{demonstram} em cada momento, favorecendo a formação teórica de quem educa.
Os \textbf{registros} constituem a memória do trabalho e a memória da instituição e valorizam as aprendizagens e o desenvolvimento dos \textbf{bebês} e das \textbf{crianças}.
Por meio deles, é possível compartilhar com as famílias e com a comunidade as experiências desenvolvidas na instituição, o que valoriza o trabalho pedagógico.
Pelo \textbf{registro} ainda é possível \textbf{documentar} os fazeres, explicitando os porquês das ações ; discutir com os demais professores aquilo que se observa como forma de ampliar conhecimentos e produzir teorias, revisitando as práticas.
Com a avaliação, proporcionamos também aos sujeitos do processo a possibilidade de rever e de pensar sobre suas experiências, compreendendo melhor o que fazem.
A documentação ainda permite que o trabalho pedagógico se individualize, atendendo a cada \textbf{bebê} e \textbf{criança} em suas necessidades específicas.
É importante compreender que a documentação pedagógica é um instrumento de avaliação.
Ela pode ser feita de diferentes maneiras : por intermédio das anotações sobre os principais fatos, falas das \textbf{crianças}, elementos interessantes de uma experiência proposta em um \textbf{diário} ; usando fotografias ; filmagens e por meio das produções dos \textbf{bebês} e \textbf{crianças}, como desenhos, pesquisas, cartazes, \textbf{pinturas}, modelagens etc.
Os \textbf{registros} do que acontece não são a base para avaliar os \textbf{bebês} e \textbf{crianças}, no sentido de dizer o que ela sabe ou não sabe, mas para pensar como o ambiente pedagógico, os espaços, o tempo, os materiais utilizados e as experiências propostas podem ser pensadas para ampliar a aprendizagem e o desenvolvimento, de acordo com os objetivos elencados para o trabalho pedagógico.
A avaliação na Educação \textbf{Infantil} não visa classificar os \textbf{bebês} e \textbf{crianças} entre as que sabem mais e as que não sabem.
Também não tem como objetivo traçar perfis ( a \textbf{criança} é agressiva ; ou a \textbf{criança} é atenciosa ; ou é desobediente, por exemplo ).
Após a aprovação da Lei n.º 12.796, de 4 de abril de 2013, alterou -se o texto da LDBEN n.º 9.394/96, Art. 31 : > A Educação \textbf{Infantil} será organizada de acordo com as seguintes regras comuns : > I – avaliação mediante \textbf{acompanhamento} e \textbf{registro} do desenvolvimento das \textbf{crianças}, sem o objetivo de promoção, mesmo para o acesso ao ensino fundamental ; > II – carga horária \textbf{mínima} anual de 800 ( oitocentas ) horas, distribuída por um mínimo de 200 ( duzentos ) dias de trabalho educacional ; > III – atendimento à \textbf{criança} de, no mínimo, 4 ( quatro ) horas \textbf{diárias} para o turno parcial e de 7 ( sete ) horas para a jornada integral ; > IV – controle de frequência pela instituição de educação \textbf{pré-escolar}, exigida a frequência \textbf{mínima} de 60 \% ( sessenta por cento ) do total de horas ; > V – expedição de documentação que permita atestar os processos de desenvolvimento e aprendizagem da \textbf{criança}.
A legislação reforça o cumprimento do calendário letivo para a Educação Básica e enfatiza o lugar da avaliação na Educação \textbf{Infantil}, que não deve ser usado para a promoção para o Ensino Fundamental.
Confirma a carga horária \textbf{mínima} prevista nas DCNEIs ( BRASIL, 2009b ), aponta sobre o controle da frequência e a necessidade de documentação de \textbf{registro} do desenvolvimento da \textbf{criança}.
O processo de aprendizagem e desenvolvimento está acontecendo ininterruptamente e as experiências propostas, a forma como os espaços, tempo e materiais são organizados e a forma de como os \textbf{adultos} lidam afetivamente com os \textbf{bebês} e as \textbf{crianças} tem grande influência sobre como eles se relacionam com aquilo que lhes é proposto nas \textbf{creches} e \textbf{pré-escolas}.
Avaliar o planejamento e a prática pedagógica, assim como as condições físicas, de estrutura e de pessoal é condição para termos instituições cada vez mais promotoras dos direitos de aprendizagem e desenvolvimento que deve ser garantido aos \textbf{bebês} e \textbf{crianças} por meio dos objetivos dos campos de experiências.
Nossas intervenções devem ter como foco a proposição de novas experiências quando percebemos que elas não estão mais instigando e desafiando os \textbf{bebês} e \textbf{crianças}.
É preciso considerar : - Os materiais : quando se percebe que as \textbf{crianças} não se \textbf{interessam} mais por manipulá -los ou \textbf{brincar} com eles, é indicativo que já os exploraram de todas as formas ; - O tempo : quando se percebe que estes são longos ou curtos demais para o desenvolvimento das atividades com as quais os \textbf{bebês} e as \textbf{crianças} são convidados a participar ; - As formas de explorar as experiências : quando se percebe que estas não trazem um apelo à \textbf{curiosidade} e não envolvem afetivamente os \textbf{bebês} e \textbf{crianças} ; - Os espaços : quando estão marcados pela imobilidade, monotonia e não incentivam os \textbf{bebês} e as \textbf{crianças} a reagirem de formas diferentes aos seus estímulos ; - As relações entre \textbf{crianças} e \textbf{adultos} : quando se percebe que estão pautadas pela ausência de diálogo e mútua colaboração ; - As relações entre as \textbf{crianças} : quando percebemos que não estão sendo mediadas por novas atividades que permitam agir de maneiras novas, mais coletivas e colaborativas.
É a partir de uma avaliação de contexto que podemos pensar outros modos de atuar, outras informações a explorar, outras atividades que, modificadas, possam se constituir como estímulos para outras relações dos \textbf{bebês} e \textbf{crianças} com as experiências vividas na instituição.

% matched lemmas: acompanhamento, adulto, bebê, brincar, creche, criança, curiosidade, demonstrar, diário, documentar, infantil, interessar, pequeno, pintura, pré-escola, pré-escolar, registro, rotina
\end{textsample}
